\setcounter{page}{1}
\chapter{Introduzione}

Il progetto, realizzato dal nostro gruppo Error404 (Lorenzo Sabatino, Damiano Pronti, Guido Francesco Valente, Luca Picano, Patrik Perkowski), si concentra sul tema delle infrastrutture verdi, scelto dopo un’analisi iniziale di tre ambiti proposti: prodotti bio, infrastrutture verdi e mobilità sostenibile.
Le infrastrutture verdi comprendono spazi come parchi, giardini, piste ciclabili e aree pedonali, che combinano natura e urbanizzazione per migliorare la qualità della vita e promuovere la sostenibilità. Questi spazi, fondamentali per la green economy, aiutano a:
\begin{itemize}
    \item Ridurre l’impatto ambientale delle città.
    \item Migliorare la qualità dell’aria.
    \item Incentivare attività sostenibili e la valorizzazione del verde.
\end{itemize}
Abbiamo scelto questo tema per via delle sue criticità, come la difficoltà nel reperire informazioni aggiornate o trovare nuove infrastrutture verdi e un maggior riscontro di problemi rispetto ai temi dei prodotti bio e della mobilità sostenibile. La scelta è stata inoltre   supportata da:
\begin{itemize}
    \item Ampia libertà di sviluppo.
    \item Accessibilità degli utenti intervistati.
    \item Risultati positivi emersi dalle interviste preliminari.
\end{itemize}
Abbiamo svolto interviste preliminari con domande aperte, elencate di seguito, volte ad esplorare l’esperienza delle persone con le infrastrutture verdi; le domande erano mirate a raccogliere informazioni generali e dettagli sull’uso di questi spazi, includendo aspetti come la ricerca, l’accesso e il livello di soddisfazione.
Domande preliminari poste:
\begin{enumerate}
    \item Come ti chiami?
    \item Quanti anni hai?
    \item Da dove vieni?
    \item Quando è stata l’ultima volta che hai utilizzato un’infrastruttura verde?
    \item Raccontami la tua ultima esperienza in un’infrastruttura verde.
    \item Come l’hai trovata?
    \item Come hai raggiunto il posto?
    \item Perché l’hai utilizzata?
    \item Su una scala da 1 a 5, come valuti la tua esperienza? Perché?
\end{enumerate}
Le risposte ottenute ci hanno permesso di delineare il profilo degli utenti ai quali eravamo interessati, per il progetto, suddiviso in quattro principali categorie:
\begin{itemize}
    \item Cittadini urbani, che utilizzano gli spazi verdi per attività sportive e ricreative.
    \item Famiglie, alla ricerca di aree sicure e attrezzate per i bambini.
    \item Turisti, interessati a scoprire nuovi luoghi naturali in città.
    \item Comunità locali e gruppi, che organizzano eventi o iniziative collettive in spazi verdi.
    \item Sportivo, con la voglia di svolgere attività all’aperto.
\end{itemize}
Per suddividere nel migliore dei modi il lavoro di gruppo e per rispettare le scadenze  abbiamo creato ed utilizzato una timeline per l’intero progetto.
